% Options for packages loaded elsewhere
\PassOptionsToPackage{unicode}{hyperref}
\PassOptionsToPackage{hyphens}{url}
\PassOptionsToPackage{dvipsnames,svgnames,x11names}{xcolor}
%
\documentclass[
  11pt,
]{article}
\usepackage{amsmath,amssymb}
\usepackage{iftex}
\ifPDFTeX
  \usepackage[T1]{fontenc}
  \usepackage[utf8]{inputenc}
  \usepackage{textcomp} % provide euro and other symbols
\else % if luatex or xetex
  \usepackage{unicode-math} % this also loads fontspec
  \defaultfontfeatures{Scale=MatchLowercase}
  \defaultfontfeatures[\rmfamily]{Ligatures=TeX,Scale=1}
\fi
\usepackage{lmodern}
\ifPDFTeX\else
  % xetex/luatex font selection
\fi
% Use upquote if available, for straight quotes in verbatim environments
\IfFileExists{upquote.sty}{\usepackage{upquote}}{}
\IfFileExists{microtype.sty}{% use microtype if available
  \usepackage[]{microtype}
  \UseMicrotypeSet[protrusion]{basicmath} % disable protrusion for tt fonts
}{}
\makeatletter
\@ifundefined{KOMAClassName}{% if non-KOMA class
  \IfFileExists{parskip.sty}{%
    \usepackage{parskip}
  }{% else
    \setlength{\parindent}{0pt}
    \setlength{\parskip}{6pt plus 2pt minus 1pt}}
}{% if KOMA class
  \KOMAoptions{parskip=half}}
\makeatother
\usepackage{xcolor}
\usepackage[margin=1in]{geometry}
\setlength{\emergencystretch}{3em} % prevent overfull lines
\providecommand{\tightlist}{%
  \setlength{\itemsep}{0pt}\setlength{\parskip}{0pt}}
\setcounter{secnumdepth}{-\maxdimen} % remove section numbering
\ifLuaTeX
  \usepackage{selnolig}  % disable illegal ligatures
\fi
\usepackage{bookmark}
\IfFileExists{xurl.sty}{\usepackage{xurl}}{} % add URL line breaks if available
\urlstyle{same}
\hypersetup{
  pdftitle={Ordinary smooth regularization on reduced algebraic curves},
  pdfauthor={Local research continuation for P02.4},
  colorlinks=true,
  linkcolor={Maroon},
  filecolor={Maroon},
  citecolor={Blue},
  urlcolor={Blue},
  pdfcreator={LaTeX via pandoc}}

\title{Ordinary smooth regularization on reduced algebraic curves}
\author{Local research continuation for P02.4}
\date{4 October 2026}

\begin{document}
\maketitle

\section{Abstract}\label{abstract}

We give a fixed-neighborhood construction of decreasing ordinary
Chambert--Loir--Ducros smooth plurisubharmonic nets for arbitrary
classical plurisubharmonic functions on reduced proper algebraic curves
over an algebraically closed, complete, nontrivially valued field. The
construction allows several normalization branches and values equal to
minus infinity. It uses the explicitly stated finite-tree smoothing
theorem imported from the local manuscript. We prove the new
directed-selection and finite-branch descent arguments here. On a
compact auxiliary domain, strict sublevels being countable unions of
closed sets is an exact functionwise condition for a decreasing
continuous sequence. In particular, countable density of the field gives
ordinary smooth decreasing sequences in the displayed curve scope. A
saturated projective-bundle application is also proved. This is an
unreviewed local paper with an explicit imported foundation; it does not
settle the original classification of higher-dimensional analytic spaces
in P02.4.

\section{1. Domain, positivity and the imported
foundation}\label{domain-positivity-and-the-imported-foundation}

Throughout, \(K\) is algebraically closed, complete and nontrivially
valued. An ordinary open means an open subset for the Berkovich
topology. A classical psh function has values in
\(\mathbb R\cup\{-\infty\}\) and is upper semicontinuous, with the
classical curve-test property after analytic scalar extension. The
permitted identically minus-infinity case is handled explicitly. We
never replace an arbitrary such target by a finite continuous function.

An ordinary smooth function is smooth for analytic moment maps in the
Chambert--Loir--Ducros sense. Weak smoothness for harmonic
tropicalizations is a different notion. Our positive approximants are
ordinary smooth and have positive smooth curvature forms locally;
consequently they are classically psh. We use pullback stability of this
actual ordinary smooth positive-form class. At a singular point our
descended approximants will be constant on an ordinary neighborhood,
with curvature zero there.

We use the following imported statement, with precisely its hypotheses.

\textbf{Imported finite-tree smoothing theorem.} Let \(C\) be a smooth
proper algebraic \(K\)-curve. Let \(Y\subset C^{\mathrm{an}}\) be a
connected strict affinoid domain with a strictly semistable skeleton
that is a finite tree. If \(f\in C(Y)\) is lisse and
Thuillier-subharmonic on \(Y^\circ\), then, for every \(\varepsilon>0\),
there is \(g\in C(Y)\) with \[
\|g-f\|_Y<\varepsilon
\] such that \(g|_{Y^\circ}\) is ordinary smooth and psh. Each point of
an ordinary open in a smooth algebraic curve has a strict affinoid tree
neighborhood of this type contained in that open.

This is the exact local manuscript statement
\texttt{prop:p063-tree-lisse-ordinary-regularization} and its
accompanying local-tree lemma. Its statement and complete proof are
separately pinned in the source binding of the preceding fixed-domain
working proof {[}4{]}. The imported proof strictifies the finite graph
with positive rational interior masses, moves positive residual divisors
into unused directions by Riemann--Roch, realizes the resulting
principal residual divisor by one ordinary smooth harmonic germ, caps
its negative slopes on the whole unused disc component, smooths type-3
joins, and checks open overlaps. The affinoid \(Y\) is fixed before any
approximation index. Refined skeletons, collars and fresh directions may
vary with \(f\), while \(Y^\circ\) cannot vary. We invoke this imported
construction as a mathematical premise, rather than treating historical
review receipts as a proof. Independent review of the present paper must
retain that origin.

The other source inputs are the following published results, used within
their exact scopes: Thuillier's decreasing finite continuous lisse
regularization on smooth curves {[}1, Theorem 3.4.2{]};
Gubler--Rabinoff's classical/subharmonic comparison and finite
normalization description {[}2, Proposition 4.1.2 and Lemma 3.9{]};
classical analytic pullback and scalar stability {[}2, Proposition
4.2{]}; and the ordinary analytic moment-map calculus {[}3{]}. For
identically minus-infinity components we use constants \(-j\). None of
these inputs asserts that every harmonic function is ordinary smooth.

Here are two elementary consequences of the ordinary calculus needed
below. First, an ordinary smooth function on a smooth curve is locally
constant at a \(K\)-rational rigid point. In a moment chart it is a
smooth function of finitely many \(\log|f_\ell|\), with every \(f_\ell\)
an analytic unit. Each \(f_\ell(x)\) is nonzero. On a sufficiently small
ordinary neighborhood, \(|f_\ell-f_\ell(x)|<|f_\ell(x)|\), so the
ultrametric inequality gives \(|f_\ell|=|f_\ell(x)|\). The finite tuple,
and hence the smooth function of it, is constant there.

Second, regularized maxima preserve ordinary smooth psh functions. Fix a
smooth even convex \(\chi_\varepsilon\) satisfying \[
|t|\leq\chi_\varepsilon(t)\leq |t|+\varepsilon,
\qquad \chi_\varepsilon(t)=|t|\quad (|t|\geq\varepsilon),
\qquad |\chi_\varepsilon'|\leq1.
\] Such a function is obtained by convolving \(|t|\) with a nonnegative
even smooth probability density supported in
\((-\varepsilon,\varepsilon)\); outside that interval the first moment
is zero. Set \[
M_\varepsilon(a,b)=\frac{a+b+\chi_\varepsilon(a-b)}2.
\] It is smooth, convex and nondecreasing in each variable, equals the
maximum when \(|a-b|\geq\varepsilon\), and lies between \(\max(a,b)\)
and \(\max(a,b)+\varepsilon/2\). On the joint ordinary moment chart for
smooth psh functions \(a,b\), its curvature is the sum of their positive
curvature forms with nonnegative coefficients and the positive
semidefinite Hessian term for this convex composition. Thus it is an
actual ordinary smooth psh function. This proof uses neither a psh
minimum nor an assertion about all weakly smooth functions.

\section{2. Compact directed selection on a smooth
branch}\label{compact-directed-selection-on-a-smooth-branch}

\textbf{Lemma 2.1.} Let \(Y\) be compact Hausdorff. Let \((f_i)\) be a
decreasing directed family of finite continuous functions with infimum
\(u\). Let \(\mathcal A\subset C(Y)\) be stable under addition of real
constants and uniformly approximate each \(f_i\) to every positive
error. Then \[
\mathcal M_u=\{a\in\mathcal A:\text{ for some }\delta>0,\ a\geq u+\delta\text{ on }Y\}
\] is nonempty and downward directed, and its infimum is \(u\). Its
actual functions, indexed by reverse pointwise order, form a decreasing
net converging to \(u\) at every point. We interpret
\(-\infty+\delta=-\infty\).

\textbf{Proof.} If \(\|g-f_i\|_Y<\varepsilon\), then
\(g+2\varepsilon\in\mathcal M_u\), since \[
f_i+\varepsilon\leq g+2\varepsilon\leq f_i+3\varepsilon.
\] Given \(a,b\in\mathcal M_u\), choose \(\delta>0\) no larger than
their two available gaps. The function \(t=\min(a,b)-\delta/2\) is a
finite continuous comparison function with \(t>u\) everywhere. The open
sets \(\{f_i<t\}\) form an increasing directed cover of \(Y\).
Compactness and a dominating index give one \(i\) with \(f_i<t\)
everywhere. Choose \(\varepsilon<\delta/6\) and set \(c=g+2\varepsilon\)
for a uniform approximant to \(f_i\). Then \[
u+\varepsilon\leq c\leq f_i+3\varepsilon<\min(a,b).
\] This proves downward directedness. The minimum is used solely as a
comparison function; it is not asserted to be psh. For a point \(y\) and
a finite \(R>u(y)\), choose \(i\) with \(f_i(y)<R\) and then
\(\varepsilon\) with \(f_i(y)+3\varepsilon<R\). The shifted approximant
is below \(R\) at \(y\). This works for every finite \(R\) when
\(u(y)=-\infty\), and proves the claimed infimum. \(\square\)

\textbf{Lemma 2.2 (compact domination).} If \((q_\lambda)\) is a
decreasing directed net of finite continuous functions on a compact
space \(A\), with infimum \(u\), and \(t\in C(A)\) satisfies \(t>u\)
everywhere, then some \(q_\lambda<t\) on all of \(A\). Finitely many
additional strict pointwise inequalities compatible with the limit can
be imposed at the same time.

\textbf{Proof.} The sets \(\{q_\lambda<t\}\) are an increasing directed
open cover. Take a finite subcover and an index dominating its indices.
Each additional pointwise bound holds at some index by convergence;
dominate those finitely many indices too. \(\square\)

Apply Lemma 2.1 on a fixed affinoid tree
\(Y\subset W\subset C^{\mathrm{an}}\) with \(Y\) contained in a slightly
larger ordinary open on which \(u\) is defined. Thuillier provides the
decreasing finite lisse family on that larger open. Restrict it to the
same \(Y\) and apply the imported finite-tree smoothing theorem. Let
\(\mathcal A\) be the finite continuous functions on \(Y\) whose
restrictions to \(Y^\circ\) are ordinary smooth psh. This class is
stable under constants. Lemma 2.1 gives an actual ordinary smooth psh
decreasing net on \(Y^\circ\) converging to the arbitrary \(u\). There
is no countability, residue-genus or residue-Picard torsion assumption.
A component where \(u\) is identically minus infinity can instead be
treated by the constant sequence \(-j\).

\section{3. Finite normalization and several
branches}\label{finite-normalization-and-several-branches}

\textbf{Theorem 3.1.} Let \(D\) be a reduced proper algebraic
\(K\)-curve of pure dimension one, let \(W\subset D^{\mathrm{an}}\) be
ordinary open, and let \(u\) be an arbitrary classical psh function on
\(W\). At every \(x\in W\) there is one compact strict affinoid
neighborhood \(Y\subset W\), with \(x\in U=Y^\circ\), and a decreasing
net of finite continuous functions on \(Y\), ordinary smooth and psh on
the fixed \(U\), converging pointwise to \(u|_Y\). Multiple
normalization branches and poles at the singularity are allowed. Each
constructed net pulls back with the same order and limit after every
analytic scalar extension.

\textbf{Proof.} At a nonsingular point use Section 2 on its smooth
proper normalization component. Consider a singular point \(x\).
Algebraic closedness makes \(x\) and its finitely many normalization
preimages rigid and \(K\)-rational. The normalization is finite, and its
components are smooth proper algebraic curves. Shrink an ordinary
neighborhood \(W_0\subset W\) so that \(x\) is its only singularity. By
the local finite-normalization description, its inverse image is a
disjoint union of neighborhoods of points \(x_1,\ldots,x_r\), with the
normalization an isomorphism away from those points. We do not suppose
\(r=1\).

For each \(x_i\), choose a strict affinoid tree \(Z_i\) inside that
branch, with \(x_i\in O_i=Z_i^\circ\). Finiteness makes the
normalization closed. The set \[
W_1=W_0\setminus\nu\bigl(\nu^{-1}(W_0)\setminus\textstyle\coprod_i O_i\bigr)
\] is therefore ordinary open and contains \(x\). Choose a compact
strict affinoid neighborhood \(Y\subset W_1\) with \(x\in Y^\circ\);
good strict analytic local structure gives such a neighborhood. Set \[
A_i=\nu^{-1}(Y)\cap O_i,
\qquad A=\coprod_{i=1}^r A_i.
\] The inverse image \(\nu^{-1}(Y)\) is compact. Its finitely many
pieces \(A_i\) form a disjoint open partition because the whole inverse
image is contained in \(\coprod_i O_i\). Each piece is therefore also
closed in that compact inverse image, hence compact, and lies inside
\(O_i\). All later inequalities are checked on this one fixed compact
union. Put \(u_i=u\circ\nu\) on \(Z_i\) and \(a=u(x)=u_i(x_i)\). Section
2 gives on each \(Z_i\) a decreasing net \((s_{i,\lambda})\) of finite
continuous functions, ordinary smooth psh on \(O_i\), with infimum
\(u_i\). The construction for an identically minus-infinity branch uses
constants. Every term is constant on some ordinary neighborhood of
\(x_i\), by the moment-map observation in Section 1.

Let \(\mathcal B\) consist of finite continuous functions on \(Y\) that
are ordinary smooth psh on \(U\), and constant on an ordinary
neighborhood of \(x\) in \(U\). Define \[
\mathcal N_u=\{v\in\mathcal B:\ v\geq u+\delta\text{ on }Y
\text{ for some }\delta>0\}.
\] A sufficiently large constant belongs to this class: upper
semicontinuity bounds \(u\) above on the compact \(Y\). We prove that
\(\mathcal N_u\) is downward directed and has infimum \(u\) by actual
branch constructions.

Take \(v,w\in\mathcal N_u\) and choose \(\delta>0\) no larger than their
two gaps. Write \(t=\min(v,w)\) on \(Y\) and \(t_i=t\circ\nu\) on
\(A_i\). We have \(t_i\geq u_i+\delta\). Fix \(0<\varepsilon<\delta/8\).

If \(a\) is finite, compact domination on \(A_i\), together with
pointwise convergence at \(x_i\), gives a term \(g_i\) satisfying \[
g_i<t_i-4\varepsilon\quad\text{on }A_i,
\qquad a\leq b_i:=g_i(x_i)<a+\varepsilon.
\] Let \(b=\max_i b_i\) and put \[
h_i=g_i+(b-b_i)+\varepsilon.
\] Then \(0\leq b-b_i<\varepsilon\), each \(h_i\geq u_i+\varepsilon\),
and \[
h_i< t_i-2\varepsilon\quad\text{on }A_i.
\] Every \(h_i\) is ordinary smooth psh on \(O_i\), constant near
\(x_i\), and has the common value \(b+\varepsilon\) there.

If \(a=-\infty\), choose a finite \(c\) with \[
c<\min_A t\circ\nu-4\varepsilon.
\] Compact domination and convergence at \(x_i\) give \(g_i\) with \[
g_i<t_i-4\varepsilon\quad\text{on }A_i,
\qquad g_i(x_i)<c-\varepsilon.
\] Set \[
h_i=M_\varepsilon(g_i,c)+\varepsilon.
\] These are ordinary smooth psh functions and satisfy \[
u_i+\varepsilon\leq h_i
\leq\max(g_i,c)+\tfrac32\varepsilon
<t_i-\tfrac52\varepsilon
\quad\text{on }A_i.
\] They equal the common constant \(c+\varepsilon\) on actual
neighborhoods of every \(x_i\). This construction includes branches
identically equal to minus infinity.

In either case, the \(h_i|_{A_i}\) give a continuous function \(z\) on
\(Y\). Indeed, the finite map \(A\to Y\) is a closed quotient map, is
bijective away from \(x\), and identifies only the displayed preimages
of \(x\), where all values agree. Their constant neighborhoods upstairs
descend to an actual ordinary constant neighborhood downstairs: remove
the image of the closed complement of those neighborhoods under the
finite normalization. Intersecting with \(U\) gives a neighborhood of
\(x\) on which \(z\) is constant. Away from \(x\) use the normalization
isomorphism. These ordinary charts agree on overlaps, so \(z\) is
ordinary smooth psh on \(U\), with curvature zero near \(x\) and
positive smooth curvature elsewhere. Thus \(z\in\mathcal N_u\) and
\(z<\min(v,w)\) on \(Y\). This proves downward directedness without
asserting general smooth descent along an arbitrary finite map.

It remains to check the infimum at every point. Start with a fixed large
constant in \(\mathcal N_u\), use it for both \(v,w\) above, and fix
\(y\in Y\) and a finite \(R>u(y)\). If \(a\) is finite and \(y=x\),
choose \(\varepsilon\) so small that \(a+2\varepsilon<R\). The
constructed common value is \(b+\varepsilon<a+2\varepsilon<R\). If \(a\)
is finite and \(y\ne x\), the point \(y\) has one normalization preimage
\(\widetilde y\). Choose \(\varepsilon\) with \(u(y)+2\varepsilon<R\)
when \(u(y)\) is finite, or any sufficiently small positive
\(\varepsilon\) when \(u(y)=-\infty\). Impose also
\(g_i(\widetilde y)<R-2\varepsilon\) on its branch. Lemma 2.2 permits
that simultaneous constraint and gives \(h_i(\widetilde y)<R\).

If \(a=-\infty\), choose in addition \(c<R-2\varepsilon\). For \(y=x\)
the descended value is \(c+\varepsilon<R\). For \(y\ne x\) impose
\(g_i(\widetilde y)<R-2\varepsilon\) as well, taking \(\varepsilon\) so
that this is compatible with its limit. The maximum bound then gives
\(h_i(\widetilde y)<R\). All choices involve finitely many branches and
finitely many pointwise conditions. Consequently \(\inf\mathcal N_u=u\)
on all of \(Y\), including poles. Indexing \(\mathcal N_u\) by reverse
pointwise order proves the theorem's net assertion.

The functions are actual local functions on \(U\). Under an analytic
morphism, analytic moment maps and their positive curvature forms pull
back; constants remain constant. Their restrictions to overlapping
charts remain equal. In particular the same functions pull back to
\(U_L\) for each analytic extension \(L/K\), and pointwise order and
convergence to \(u\circ p\) hold at each point of \(U_L\) by evaluation
at its image. This is the pullback of the constructed base-field net,
rather than descent of arbitrarily chosen upstairs approximants.
\(\square\)

\section{4. Exactly when a sequence can be
extracted}\label{exactly-when-a-sequence-can-be-extracted}

\textbf{Lemma 4.1.} Let \(Y\) be compact Hausdorff and let
\((q_\lambda)\) be any decreasing directed net of finite continuous
functions with infimum \(u\). There is a decreasing sequence of terms of
this net converging pointwise to \(u\) if every rational strict sublevel
\(\{u<r\}\) is \(F_\sigma\) in \(Y\). Conversely, the existence of any
pointwise decreasing continuous sequence with limit \(u\) implies that
condition.

\textbf{Proof.} Necessity follows from \[
\{u<r\}=\bigcup_j\{q_j<r\}
=\bigcup_{j,k\geq1}\{q_j\leq r-1/k\},
\] whose last sets are closed. For sufficiency, write
\(\{u<r\}=\bigcup_k F_{r,k}\) with closed, hence compact, \(F_{r,k}\).
Each nonempty \(F_{r,k}\) is covered by the increasing directed open
sets \(\{q_\lambda<r\}\). A finite subcover and a dominating index yield
one index \(\lambda(r,k)\) below \(r\) on that whole closed piece.
Enumerate these countably many indices, include an arbitrary initial
index, and recursively take dominating indices. The resulting sequence
decreases. If its infimum at a point exceeded \(u\), an intermediate
rational \(r\) and a closed piece containing that point would contradict
the eventual domination of \(\lambda(r,k)\). The same argument works for
\(u=-\infty\) using every rational \(r\). \(\square\)

Apply Lemma 4.1 to the actual net of Theorem 3.1 on its compact \(Y\).
Its terms are already ordinary smooth psh on the fixed \(U\); selecting
countably many dominating indices does not change that domain or destroy
branch alignment. Thus the \(F_\sigma\) condition is an exact
functionwise criterion for a decreasing sequence of finite continuous
functions on this \(Y\), ordinary smooth psh on \(U\), with limit \(u\)
on all of \(Y\).

If \(K\) has a countable dense subset, every strict affinoid \(Y\) is
compact metrizable. To see this, a countable family of Tate-polynomial
evaluations densely approximates the elements of its affinoid algebra
and separates bounded seminorms. Those evaluations give an embedding
into a countable product of compact intervals; compactness and
Hausdorffness give the stated metrizable topology. Strict sublevels of
the upper semicontinuous \(u\) are open, and every open subset of a
metric space is \(F_\sigma\) (use closed positive-distance inner
layers). We obtain the following consequence.

\textbf{Corollary 4.2.} If \(K\) is also countably dense, every
arbitrary classical psh function on an ordinary open in a reduced proper
algebraic \(K\)-curve has, locally on one fixed ordinary neighborhood, a
pointwise decreasing sequence of finite ordinary smooth psh
approximants. This includes every finite normalization branch count and
every permitted minus-infinity target value. If the target is finite
continuous on the compact auxiliary \(Y\), convergence there is uniform
by Dini's theorem.

The compact-domain qualifier matters. A sequence defined only on \(U\)
need not continuously extend to this particular \(Y\). Necessity is
therefore asserted for sequences on \(Y\), or on a neighborhood of a
smaller compact domain \(Y'\subset U\), rather than for an arbitrary
unextended sequence on \(U\). No topology-only criterion for all
higher-dimensional sheaf germs follows.

\section{5. A saturated higher-dimensional
application}\label{a-saturated-higher-dimensional-application}

Let \(p:X\to D^{\mathrm{an}}\) be an analytically locally trivial
projective-space bundle of relative dimension \(m\geq1\). Let
\(V\subset D^{\mathrm{an}}\) be ordinary open and let \(v\) be an
arbitrary classical psh function on the entire saturated domain
\(p^{-1}(V)\). At a base point \(c\), its restriction to the fiber is
classical psh on \(\mathbb P^m_{\mathcal H(c)}\). The complete residue
field \(\mathcal H(c)\) need not be algebraically closed. That fiber is
compact, connected and boundaryless. If the restriction is not
identically minus infinity, upper semicontinuity gives a finite attained
maximum. By the strong maximum principle {[}2, Theorem 6.1{]}, its
maximum locus is open; it is also closed and nonempty. Connectedness
makes the restriction constant. This proves constancy on every fiber,
including fibers over nonrigid base points.

There is consequently one function \(u\) on \(V\) with \(v=u\circ p\).
Analytic local sections show that \(u\) is upper semicontinuous and
classically psh, and its sections agree on overlaps by fiber constancy.
Apply Theorem 3.1 at a base point and pull back its ordinary smooth psh
net to \(p^{-1}(U)\). Under countable density apply Corollary 4.2 to
obtain a sequence. Every term has the pulled-back actual positive smooth
curvature form, retains order and pointwise limit, and behaves likewise
after scalar extension.

This proves a theorem for the full classical class on saturated domains
of these bundles, without any torus invariance assumption. An arbitrary
psh germ on a nonsaturated affine fiber chart can vary along the fiber,
so this argument makes no claim about that class.

\section{6. Scope and remaining original
obligations}\label{scope-and-remaining-original-obligations}

The imported finite-tree smoothing is an explicit foundation from the
authorized local manuscript. The directed-selection, functionwise
extraction and multibranch descent above are new local proof
applications. No independent verdict is claimed for this paper, and no
frozen native draft, source or canonical note has been altered. Neither
active native equation author received this paper or its imported
premise.

P02.4 asks for the sharp geometric range of ordinary smooth decreasing
psh regularization starting with every classical psh function. The
present theorem covers reduced proper algebraic curves over the
displayed algebraically closed field and a saturated bundle sector. It
does not characterize nonalgebraizable curve ranges, remove algebraic
closedness in the ordinary construction, classify general
higher-dimensional local germs, or settle analytic-boundary cases. It
therefore does not close P02.4.

\section{References and exact origin}\label{references-and-exact-origin}

\begin{enumerate}
\def\labelenumi{\arabic{enumi}.}
\tightlist
\item
  A. Thuillier, \emph{Théorie du potentiel sur les courbes en géométrie
  analytique non archimédienne}, 2005, HAL \texttt{tel-00010990v1}.
  Theorem 3.4.2 is used for a net on a fixed ordinary open, not silently
  for a sequence.
\item
  W. Gubler and J. Rabinoff, arXiv \texttt{2506.13548v2}, cached primary
  PDF and literal extraction. Proposition 4.1.2, Lemma 3.9, Proposition
  4.2 and Theorem 6.1 are the exact cited inputs.
\item
  A. Chambert-Loir and A. Ducros, arXiv \texttt{1204.6277v2}, saved
  revision of 25 July 2025. The ordinary moment-map calculus and
  pullback of positive smooth forms are used.
\item
  Imported local manuscript packet,
  \nolinkurl{p063_curve_cld_monotone_regularization_statement_20260827.tex}
  and
  \nolinkurl{p063_curve_cld_monotone_regularization_proof_20260827.tex}.
  The exact statement SHA-256 is
  \nolinkurl{767cbf81c723279c1218df4136f8e299465240b325f41f6b0663d001fb5d8a68};
  the complete proof SHA-256 is
  \nolinkurl{d1db3c2db861d920598697fdad5d002a2b1ce947dbdbd95ad9b78b8f9f312bd8}.
  The earlier local application is preserved in
  \nolinkurl{P02.4-fixed-domain-net-and-sequence-bridge-20261004.md} and
  its source binding.
\end{enumerate}

\end{document}
